\bookchapter{บทที่ 10: การชำระเงินผ่านมือถือกับบริการชีวิตประจำวัน}

\chapterintro{พูดตรงๆ รุ่นพ่อรุ่นแม่ฉันต่อต้านเรื่อง “จ่ายเงินผ่านมือถือ” ตั้งแต่แรก “เงินอยู่ในมือถือ ถ้าโดนขโมยจะทำยังไง?” ความกังวลนี้เป็นเรื่องปกติ แต่ความจริงคือ: เงินสดในกระเป๋ากางเกงคุณโดนหายง่ายกว่าการที่มือถือโดนแอบรูดมากๆ การชำระเงินผ่านมือถือไม่ได้เอาเงินไปเก็บในมือถือ มันเป็นแค่กุญแจดอกหนึ่ง กุญแจที่เปิดบัญชีธนาคารของคุณ}

\booksection{หนึ่ง ผูกบัตรธนาคารกันแน่ยังไง?}

ไม่ว่าคุณจะใช้ Alipay หรือ WeChat การผูกบัตรครั้งแรกต้องผ่านขั้นตอนพวกนี้:

\begin{enumerate}
  \item เปิดแอป หา "ของฉัน" หรือ "Wallet" กด "เพิ่มบัตรธนาคาร"
  \item กรอกเลขบัตร ตอนนี้แอปส่วนใหญ่รองรับถ่ายรูปจำแนกเลขบัตรได้เลย ไม่ต้องพิมพ์ทีละตัว
  \item ธนาคารจะส่ง SMS รหัสยืนยันมาที่เบอร์ที่ลงทะเบียนไว้ กรอกลงไป แล้วตั้งรหัสชำระเงิน 6 หลัก
  \item รหัสชำระเงินตัวนี้กับรหัสบัตรธนาคาร รหัสล็อกหน้าจอ ควรเป็นรหัสสามตัวที่ต่างกัน อย่าขี้เกียจตั้งเหมือนกันหมด
\end{enumerate}

มีเคล็ดลับเล็กๆ: ถ้ากังวลเรื่องความปลอดภัย ไปเปิดบัญชีอิเล็กทรอนิกส์ประเภท II หรือ III ในแอปธนาคาร ไว้ผูกการชำระเงินผ่านมือถือโดยเฉพาะ ใส่แค่เงินใช้ประจำวัน แม้เกิดเรื่องขึ้น ความเสียหายก็จำกัด

\booksection{สอง คิวอาร์จ่ายเงินคือ “ใครเห็นใครรับเงิน”}

หลายคนแยกไม่ออกระหว่างคิวอาร์จ่ายกับคิวอาร์รับ ความต่างนี้สำคัญมาก

\begin{itemize}
  \item \textbf{คิวอาร์จ่ายเงิน}: บาร์โค้ด + คิวอาร์ที่คุณโชว์ให้แคชเชียร์สแกน คิวอาร์ตัวนี้เทียบเท่า "กระเป๋าเงินดิจิทัล" ของคุณ ใครสแกนก็หักเงินจากบัญชีคุณได้ ดังนั้น\textbf{ห้ามสกรีนช็อตส่งให้ใครเด็ดขาด}. คิวอาร์ตัวนี้รีเฟรชเองทุกนาที ก็เพื่อกันการถูกนำไปใช้โดยไม่ได้รับอนุญาต
  \item \textbf{คิวอาร์รับเงิน}: คนอื่นสแกนตัวนี้เพื่อโอนเงินให้คุณ ตัวนี้สกรีนช็อตเก็บไว้ส่งให้ใครก็ได้ ปลอดภัย
\end{itemize}

จำความต่างนี้ให้ขึ้นใจ ช่วยเลี่ยงกับดักกลโกงได้เยอะ ใครก็ตามที่ให้คุณ “ส่งสกรีนช็อตคิวอาร์จ่ายเงินมา” ร้อยทั้งร้อยคือมิจฉาชีพ

\booksection{สาม การตั้งค่าความปลอดภัยเหล่านี้ ใช้เวลา 5 นาทีตั้งแล้วสบายใจขึ้นเยอะ}

\begin{itemize}
  \item \textbf{ชำระเงินด้วยลายนิ้ว/ใบหน้า}: ในการตั้งค่าการชำระเงินของ Alipay กับ WeChat เปิดการชำระเงินด้วยไบโอเมตริกซ์ ต่อจากนั้นทุกครั้งที่จ่ายต้องกดลายนิ้วหรือสแกนหน้า แม้มือถือหาย คนเก็บได้ก็ใช้เงินคุณไม่ได้
  \item \textbf{กุญแจความปลอดภัย}: ทั้ง Alipay และ WeChat มีฟังก์ชัน "กุญแจความปลอดภัย" เปิดแล้วเข้าหน้า "Wallet" ของแอปต้องยืนยันตัวตนอีกรอบ เส้นทาง: Alipay ของฉัน ตั้งค่า บัญชีและความปลอดภัย ศูนย์ความปลอดภัย กุญแจความปลอดภัย
  \item \textbf{ตรวจการชำระเงินไร้รหัสผ่าน}: หลายคนเซ็นรับบริการชำระเงินอัตโนมัติไปมากมายจนลืม ไปที่ "การตั้งค่าการชำระเงิน" ของ Alipay และ WeChat หา "การชำระเงินไร้รหัสผ่าน/หักบัญชีอัตโนมัติ" ดูว่าเซ็นบริการอะไรไว้บ้าง ปิดตัวที่ไม่จำเป็นทิ้ง
  \item \textbf{การจัดการอุปกรณ์ที่ล็อกอิน}: ดูเป็นระยะว่าอุปกรณ์ไหนเคยล็อกอินบัญชีคุณ ถ้าพบอุปกรณ์แปลกหน้า ไล่ออกทันที เส้นทาง: Alipay ตั้งค่า บัญชีและความปลอดภัย การจัดการอุปกรณ์ที่ล็อกอิน
\end{itemize}

\booksection{สี่ รถเมล์รถไฟใต้ดิน: มือถือก็คือบัตรโดยสาร}

ตอนนี้ออกไปนั่งรถเมล์รถไฟใต้ดิน แทบไม่มีใครใช้บัตรโดยสารแบบแข็งแล้ว แตะประตูกั้นด้วยมือถือมีสองเส้นทาง:

\begin{itemize}
  \item \textbf{บัตรโดยสารเสมือนใน Wallet มือถือ}: เปิดใช้บัตรโดยสารเสมือนได้เลยใน Apple Wallet, Huawei Wallet, Xiaomi Wallet จุดเด่นของสูตรนี้: แตะโดยไม่ต้องปลดล็อกมือถือ เอาเครื่องชนไปได้เลย ตอบสนองเร็วที่สุด และไม่มีเน็ตก็แตะได้ บัตรโดยสารเสมือนบางเมืองยังมีส่วนลดเปลี่ยนขบวนเหมือนบัตรจริง
  \item \textbf{คิวอาร์โดยสารของ Alipay/WeChat}: ค้นหา "คิวอาร์โดยสาร" ในแอป เปิดใช้คิวอาร์รถเมล์หรือรถไฟใต้ดินของเมืองคุณ จุดเด่น: ไม่ต้องเติมเงิน หักจากยอดคงเหลือหรือบัตรธนาคารที่ผูกไว้โดยตรง แต่ต้องต่อเน็ต และก่อนแตะต้องเปิดแอปหาคิวอาร์ก่อน
\end{itemize}

ฉันแนะนำบัตรโดยสารเสมือนใน Wallet มือถือมากกว่า แม้ครั้งแรกต้องเติมสัก 10-20 หยวน แต่ใช้จริงสะดวกกว่าเยอะ โดยเฉพาะตอนวิ่งไปขึ้นรถไฟใต้ดิน ไม่ต้องหยุดปลดล็อกมือถือหาแอป

\booksection{ห้า ค่าใช้จ่ายชีวิต สั่งอาหาร เรียกรถ}

\begin{itemize}
  \item \textbf{จ่ายค่าน้ำค่าไฟค่าแก๊ส}: ใน "ค่าใช้จ่ายชีวิต" ของ Alipay หรือ WeChat ผูกเลขบัญชีน้ำ ไฟ แก๊สของบ้านให้เรียบร้อย ต่อจากนั้นทุกเดือนได้บิลมากดเข้าไปจ่ายได้เลย เลขบัญชีมักเจอบนบิลกระดาษของบ้าน หรือโทรถามเจ้าหน้าที่
  \item \textbf{เทียบราคาสั่งอาหาร}: ร้านเดียวกันบน Meituan กับ Ele.me ราคากับโปรโมชันมักไม่เท่ากัน ก่อนสั่งดูทั้งสองฝั่ง บางทีประหยัดได้หลายหยวน อีกอย่างฟังก์ชัน "Pinhaofan" ของ Meituan ชวนเพื่อนสั่งพร้อมกัน ราคาถูกลงได้เยอะ
  \item \textbf{ความปลอดภัยเรียกรถ}: ใน Didi หรือ Amap อย่าลืมตั้ง "ผู้ติดต่อฉุกเฉิน" และ "แชร์เส้นทางเดินทาง" ไว้ ตอนเรียกรถแอปจะแชร์เส้นทางให้ผู้ติดต่อที่ตั้งไว้อัตโนมัติ อีกอย่างตอนดึกให้เลือก "แท็กซี่" ก่อน "รถด่วน" แท็กซี่มีหน่วยงานกำกับดูแล ปลอดภัยกว่าหน่อย
\end{itemize}

\booksection{หก ใช้มือถือจัดการประกันสุขภาพ ประกันสังคม และราชการ}

สองสามปีนี้เปลี่ยนไปมาก กิจกรรมที่เมื่อก่อนต้องไปต่อแถวที่หอประชุม ตอนนี้ใช้มือถือสองสามนาทีจบ

\begin{itemize}
  \item \textbf{บัตรประกันสุขภาพอิเล็กทรอนิกส์}: ค้นหา "บัตรประกันสุขภาพอิเล็กทรอนิกส์" ใน Alipay หรือ WeChat เปิดใช้งานแล้วไปหาหมอ ซื้อยาที่ร้านขายยา แตะมือถือได้เลย ไม่ต้องพกบัตรจริง ครั้งแรกต้องยืนยันใบหน้า หลังจากนั้นทุกครั้งที่ใช้จะสร้างรหัสไดนามิกให้
  \item \textbf{ตรวจประกันสังคม}: ดาวน์โหลดแอป "Zhangshang 12333" หรือใช้ "ศูนย์พลเมือง" ใน Alipay ตรวจประวัติการจ่ายเงินประกันสังคม ยอดเงินบัญชีบำนาญ การขอเบี้ยว่างงาน ฯลฯ ได้
  \item \textbf{ภาษีเงินได้}: การหักภาษีประจำปีระหว่างมีนาคมถึงมิถุนายน จัดการได้ในแอป "ภาษีเงินได้บุคคลธรรมดา" อย่าลืมดูว่ามีสิทธิ์หักเพิ่มเติมที่ประกาศได้ไหม (การศึกษาบุตร ดอกเบี้ยบ้าน เลี้ยงดูผู้สูงอายุ ฯลฯ) พวกนี้ช่วยให้จ่ายภาษีน้อยลงหรือได้คืน
\end{itemize}

\begin{tcbwarn}{จำข้อพวกนี้ไว้ กันโกงให้ปลอดภัย}
\begin{itemize}
  \item ตำรวจ อัยการ ศาล ธนาคาร เจ้าหน้าที่ จะไม่ให้คุณโอนเงินไป "บัญชีที่ปลอดภัย" เด็ดขาด
  \item ใครก็ตามที่ให้คุณส่งสกรีนช็อตคิวอาร์จ่ายเงิน คือมิจฉาชีพ
  \item อย่ากดลิงก์แปลกปลอม โดยเฉพาะ SMS แบบ "บัตร ETC หมดอายุ" "ระงับการใช้บัตรประกันสังคม"
  \item เงินคืนจากช้อปปิ้งจะคืนทางเดิม ไม่ต้องให้เลขบัตรธนาคารกับใคร
\end{itemize}

\end{tcbwarn}

\begin{tcbpractice}{ทำเลยตอนนี้: ตรวจสุขภาพความปลอดภัยให้แอปชำระเงินของคุณ}
\begin{enumerate}
  \item เปิด Alipay กับ WeChat ยืนยันว่าเปิดการชำระเงินด้วยลายนิ้ว/ใบหน้าแล้ว
  \item หาฟังก์ชัน "กุญแจความปลอดภัย" แล้วเปิดมัน
  \item ไล่ดูรายการ "การชำระเงินไร้รหัสผ่าน/หักบัญชีอัตโนมัติ" ปิดบริการที่ไม่ใช้แล้ว
  \item ดูรายการ "อุปกรณ์ที่ล็อกอิน" ถอดอุปกรณ์ที่ไม่รู้จักออก
  \item เปิดใช้งานบัตรประกันสุขภาพอิเล็กทรอนิกส์ของคุณ ครั้งหน้าไปโรงพยาบาลไม่ต้องพกบัตรจริงแล้ว
\end{enumerate}

\end{tcbpractice}
